\documentclass[11pt]{article}
\usepackage[a4paper,margin=23mm]{geometry}
\usepackage{amsmath,amssymb,amsthm,mathtools,bm,booktabs,microtype,xurl,hyperref}
\hypersetup{colorlinks=true,linkcolor=blue,urlcolor=blue,citecolor=blue,
pdftitle={P54 common renormalization, radial obstruction and CHN closure}}
\setlength{\emergencystretch}{2em}
\newcommand{\tr}{\operatorname{tr}}
\newcommand{\diag}{\operatorname{diag}}
\newcommand{\MS}{\overline{\mathrm{MS}}}
\newcommand{\dd}{\mathrm d}
\newcommand{\LOOP}{\mathcal L}
\newtheorem{proposition}{Proposition}
\newtheorem{theorem}{Theorem}
\title{Route-F P54: Common Renormalization,\\
a Radial Stability Obstruction and $C_{HN}$ Threshold Closure}
\author{P54PQ-v2 four-dimensional mainline}
\date{12 September 2026}
\begin{document}
\maketitle
\begin{abstract}
We verify the finite-parameter conversion linking the existing fixed-VEV
and minimal-subtraction descriptions of the same bare scalar action.
Deleting the large finite counterterm while retaining the original
renormalized parameters is not such a conversion. We calculate the
previously untested three-dimensional physical radial bosonic Hessian.
It has a negative eigenvalue in the declared one-loop hard-potential
truncation. An analytical bound proves that no choice of the three
Majorana masses can remove this negative direction in that
truncation at the frozen scalar point. This is a conditional benchmark
obstruction, not an all-orders or gauge-independent exclusion of the model.
Separately, we calculate the linear-$C_{HN}$ finite Yukawa vertex and
associated terminal Weinberg and Higgs-mass contributions. At partial
sterile thresholds the field-dependent Schur map generates a dimension-six
$LNH(H^\dagger H)$ operator. Its mass-dependent mixing explains a precise
residual left by a dimension-five-only sequential treatment when retained
Majorana mass effects are not additionally expanded away.
\end{abstract}

\section{Scope and fixed input}
The action remains P54PQ-v2:
$3\,16_F+54_{H,\mathbb R}+126_{H,\mathbb C}+10_{H,\mathbb C}
+1_{H,\mathbb C}$, with the two complex symmetric family matrices
$h_{\rm raw},f_{\rm raw}$. No phenomenological fit, new family parameter,
lattice scan or replacement scalar point is introduced. The historical
broken tree vacuum and upper tree PS saddle are respectively
\begin{align}
 r_B&=(w,\sigma,v_s)=(1,0.1265,0.25)\,\omega,\\
 r_U&=(1.0008103059,0,0.2542776971)\,\omega.
\end{align}
Here the notation $w=1\,\omega$ uses the same historical reference scale
as the earlier files. Both points use the same scalar parameters.
We retain DR/$\MS$, background-field Landau gauge, and the existing
finite invariant fixed-VEV mass counterterms at
$\mu_0=0.5626028899853225\,\omega$. Write $\LOOP=16\pi^2$.

``Hard potential'' means the previously declared trace over 290 positive
scalar and 33 massive vector eigenvalues at $r_B$, with the 38 scalar
and 12 vector soft directions excluded and mixed hard/soft divided
differences retained. It is not the complete soft-IR-resummed effective
action. Its Hessian is not a pole mass. This scope applies to every
stability statement below.

\section{Finite counterterms cannot be removed at fixed numerical input}
Let $p=(\mu^2,\nu^2,\mu_s^2)$ be the three dependent mass parameters.
The canonical real-field Hessian of their finite counterterm is
\begin{equation}
 \mathsf H_{\rm CT}=\diag\left(-\Delta\mu^2 I_{54},
 -\tfrac12\Delta\nu^2 I_{252},0_{20},-\Delta\mu_s^2 I_2\right).
 \label{eq:massCT}
\end{equation}
The factor $1/2$ for the complex $126$ is required by the action's
normalization and is independently consistent with the radial and phase
Ward tests. It is not a missing factor causing the large correction.

Introduce a formal loop-counting parameter $\ell$. Holding the other
counterterm conventions fixed, equality of the bare parameters gives
\begin{equation}
 p_0=p_F+\delta p_{\rm pole}+\ell\Delta p
     =p_{\MS}+\delta p_{\rm pole}+O(\ell^2),\qquad
 p_{\MS}=p_F+\ell\Delta p+O(\ell^2).
 \label{eq:scheme}
\end{equation}
Because these mass parameters enter $V_0$ linearly, at fixed fields
\begin{equation}
 V_0(X;p_F+\ell\Delta p)
 =V_0(X;p_F)+\tfrac\ell2 X^T\mathsf H_{\rm CT}X
 \label{eq:exactmass}
\end{equation}
is an exact polynomial identity. Replacing $p_F$ by $p_{\MS}$ inside a
one-loop diagram changes it only at formal two-loop order; that order
statement does not guarantee a numerically small remainder at $\ell=1$.

\begin{center}
\begin{tabular}{lrrr}
\toprule
Parameter / $\omega^2$ & $p_F$ & $\Delta p$ & $p_F+\Delta p$\\
\midrule
$\mu^2$ & $1.4095253$ & $-0.0563307$ & $1.3531946$\\
$\nu^2$ & $0.3819020$ & $-5.4697441$ & $-5.0878421$\\
$\mu_s^2$ & $0.1719203$ & $-0.0342377$ & $0.1376826$\\
\bottomrule
\end{tabular}
\end{center}
Thus using zero finite counterterms with the old $p_F$ does not describe
the same bare point. Conversely, a literal resummation of the large
converted mass shift is not automatically a controlled implementation
of the formal scheme change. The finite scheme and VEV issues are
distinct from gauge invariance of observables \cite{vev}; the present
fixed-Landau calculation does not establish an FJ or GIVS gauge-independent
parameter construction.

\subsection{Vacuum transport and the common tadpole condition}
On the radial slice, the canonical metric and mixed parameter derivative
are
\begin{equation}
 G_r=\diag(12/5,2,1),\qquad
 A_{ri}=\frac{\partial^2V_0}{\partial r_r\partial p_i}
 =-\diag(12w/5,\sigma,v_s).
\end{equation}
Let $H_r=\partial_r^2V_0$ and $t=\partial_rV_{1,h}$ at $r_B$.
The fixed-VEV condition is exactly
\begin{equation}
 t+A\Delta p=0.
 \label{eq:tad}
\end{equation}
In the converted MS description the tree stationary point shifts by
\begin{equation}
 \delta r_{\rm tree}^{\MS}=-H_r^{-1}A\Delta p=H_r^{-1}t,
 \qquad \delta r_{\rm explicit\ loop}=-H_r^{-1}t.
 \label{eq:vevshift}
\end{equation}
Their sum vanishes at first order; neither contribution may be discarded.
The calculated tree conversion coefficient is
\begin{equation}
 \delta r_{\rm tree}^{\MS}/\omega
 =(0.306642,-25.428874,1.660472).
\end{equation}
This coefficient is far too large to assert a small shift at $\ell=1$.
Three tiny formal values $\ell=10^{-5},5\times10^{-6},2.5\times10^{-6}$
verify the stationary-point identity and quadratic remainder. They are
local scheme-identity tests, not a new physical vacuum scan.

At the upper PS saddle, only the two singlet radii $w,v_s$ carry tadpoles.
Let $H_U$ be their tree Hessian. Then the previously calculated hard-plus-CT
response separates unambiguously as
\begin{equation}
 \delta r_U=-H_U^{-1}(t_{U,h}+t_{U,\rm CT})
 =\delta r_{U,\rm loop}+\delta r_{U,\rm scheme}.
\end{equation}
The same mass counterterms apply at both backgrounds. Their upper
$\sigma=0$ tadpole cannot determine a second independent $\Delta\nu^2$.
This is the scalar-valley version of the tadpole-insertion issue in
matching \cite{matching}; that reference excludes heavy gauge bosons and
is not being used as a complete gauge-matching formula.

\subsection{Wilson matching must transform with the parameters}
For an EFT relation $q=f(p)+\ell\Delta(p)$ and finite definitions
$p'=p+\ell a$, $q'=q+\ell b$, expansion of the same relation yields
\begin{equation}
 \Delta'=\Delta+b-Df\,a.
 \label{eq:matchingconversion}
\end{equation}
The finite UV and EFT parameter cards therefore both matter. For the
upper scalar exchange $C_0=J^TD^{-1}J$, let
\begin{equation}
 D_p=\mathsf H_{{\rm CT},h}
 +\sum_{r=w,v_s}\delta r_{U,\rm scheme}^r\,\partial_rD_0.
\end{equation}
The original upper eigenbasis is held fixed as a field-coordinate chart,
so its Yukawa source tensor $J$ does not independently change. The tree
Wilson coefficient shifts by
\begin{equation}
 \delta_pC_0=-J^TD_0^{-1}D_pD_0^{-1}J,
 \qquad \Delta C_{\MS}=\Delta C_F-\delta_pC_0.
\end{equation}
Their sum is the old full first-order response. Exact-resolvent finite
differences test the tree derivative on the existing two synthetic source
cases. These identities close the stated common bosonic scheme algebra,
not the missing Wilson boxes or gauge/ghost/Yukawa diagrams.

\section{Previously untested physical radial block}
\subsection{Actual action derivatives and an independent spectral check}
Let $R$ embed the three radial coordinates in the 328 real fields and
$\mathsf H(X)=V_0''(X)$. The quartic action makes $\mathsf H$ quadratic.
Using cached first and diagonal second derivatives, three new mixed
Hessian evaluations determine all $\mathsf H_{,rs}$ exactly as a
polynomial identity. No spectrum formula is fitted.

Let $F_{c,h}$ equal $x^2[\log(x/\mu_0^2)-c]$ in neighborhoods of the
hard eigenvalues and zero near the excluded soft eigenvalues. The
spectral gap makes this a smooth local spectral function. Define
\begin{equation}
 F_{c,h}'^{[1]}(x,y)=\begin{cases}
 [F_{c,h}'(x)-F_{c,h}'(y)]/(x-y),&x\ne y,\\
 F_{c,h}''(x),&x=y.
 \end{cases}
\end{equation}
Both sums below include hard and soft indices; a mixed divided
difference need not vanish. In the common eigenbasis,
\begin{equation}
 \partial_r\partial_s\tr F_{c,h}(\mathsf H)
 =\sum_i F_{c,h}'(m_i^2)(\mathsf H_{,rs})_{ii}
 +\sum_{ij}F_{c,h}'^{[1]}(m_i^2,m_j^2)
                   (\mathsf H_{,r})_{ij}(\mathsf H_{,s})_{ji}.
 \label{eq:frechet}
\end{equation}
Scalar and vector multiplicities are $1$ and $3$, with $c=3/2$ and
$5/6$, divided by $64\pi^2$. Direct finite differences of the actual
spectral potential at two step sizes independently verify the resulting
radial Hessian, not just its symmetry.

The tree and corrected bosonic generalized squared masses solve
$H_rv=m^2G_rv$ and are
\begin{align}
 \operatorname{spec}(G_r^{-1}H_{r,0})/\omega^2
   &=(0.0135205,0.0997157,2.8291659),\\
 \operatorname{spec}(G_r^{-1}H_{r,B})/\omega^2
   &=(-0.7584480,0.0941189,3.0931051),\label{eq:radialvalues}\\
 H_{r,B}&=H_{r,0}+\Pi_{r,S}+\Pi_{r,V}+R^T\mathsf H_{\rm CT}R.
\end{align}
The actual matrix is approximately
\begin{equation}
 H_{r,B}/\omega^2=
 \begin{pmatrix}
 7.421952&-0.056046&0.057186\\
 -0.056046&-1.516518&-0.006028\\
 0.057186&-0.006028&0.094555
 \end{pmatrix}.
 \label{eq:radialmatrix}
\end{equation}
Its negative direction is orthogonal to both the gauge orbit and the PQ
phase. It is not an unremoved gauge zero mode. The relative insertion
$H_{r,0}^{-1/2}(H_{r,B}-H_{r,0})H_{r,0}^{-1/2}$ has a smallest eigenvalue
about $-56.8$. The earlier doublet-only positivity check was never a
whole-vacuum stability certificate.

The large $\Delta\nu^2$ is dominated by actual high-multiplicity scalar
clusters. Four leading degenerate groups have multiplicities
$50,44,44,36$ and contribute approximately
$-1.1353$, $-1.0839$, $-1.0669$, and $-0.7821$ to $\Delta\nu^2/\omega^2$.
These are basis-invariant spectral traces. The complete group inventory
is in the JSON report; there is no evidence of a decimal-level kernel
error explaining the instability.

\section{An all-family bound within the declared hard truncation}
At the broken radial background the actual Yukawa tensors give one
massive spinor direction per family, namely the Majorana neutrino.
The three Takagi masses are $M_i=\sigma f_i$, $f_i\geq0$.
There is no $h_{\rm raw}$ radial fermion mass, and $w,v_s$ do not enter
these masses. The scalar doublet retuning also has zero radial Hessian.
These statements are checked directly in the existing Clifford dictionary.

\begin{theorem}[Fixed-point Majorana bound]
Fix the scalar parameters, $r_B$, $\mu_0$, the hard-potential truncation
specified above, and the common fixed-VEV prescription. For any three
Majorana Takagi masses, including their zero limits, adding their
one-loop potential and associated finite tadpole counterterm leaves a
negative radial direction. This statement does not include soft-IR
completion, higher loops, pole matching, or a changed scalar point.
\end{theorem}
\begin{proof}
The fermion potential is
\begin{equation}
 V_F=-\frac1{32\pi^2}\sum_{i=1}^3M_i^4
        \left(\log\frac{M_i^2}{\mu_0^2}-\frac32\right).
\end{equation}
Differentiate with respect to $\sigma$ at fixed $f_i,\mu_0$. The finite
fixed-VEV counterterm has $\Delta\nu_F^2=(\partial_\sigma V_F)/\sigma$.
Consequently its only nonzero radial curvature is
\begin{align}
 F_{\sigma\sigma}
 &=\partial_\sigma^2V_F-\frac{\partial_\sigma V_F}{\sigma}\\
 &=-\frac4{\LOOP\sigma^2}\sum_{i=1}^3M_i^4
                        \log\frac{M_i^2}{\mu_0^2}.
 \label{eq:fermioncurvature}
\end{align}
For $x\geq0$, the continuous function $-x^2\log(x/\mu_0^2)$ has its global
maximum $\mu_0^4/(2e)$ at $x=\mu_0^2e^{-1/2}$. Thus
\begin{equation}
 F_{\sigma\sigma}\leq
 F_{\max}=\frac{6\mu_0^4}{e\LOOP\sigma^2}
          =0.08751124\,\omega^2.
 \label{eq:fermioncap}
\end{equation}
This bound allows arbitrarily large $f_i$; no perturbativity cutoff or
flavor fit is needed for the inequality. In the Loewner order,
\begin{equation}
 H_{r,B}+H_{r,F}\preceq H_{r,B}+F_{\max}e_\sigma e_\sigma^T.
\end{equation}
Already the fixed normalized direction $q=e_\sigma/\sqrt2$ obeys
\begin{equation}
 q^TG_rq=1,\qquad
 q^T(H_{r,B}+H_{r,F})q
 \leq\frac{-1.51651817+0.08751124}{2}\,\omega^2<0.
\end{equation}
An optimized generalized-eigenvector witness for the upper-bound matrix
gives the slightly stronger value $-0.7146953\,\omega^2$. This one fixed
witness works for every allowed family matrix. Hence all choices have
a negative radial direction in the stated truncation.
\end{proof}

Equation~\eqref{eq:fermioncurvature} is independently checked by direct
finite differences of $V_F$ plus its counterterm, and
Eq.~\eqref{eq:fermioncap} is saturated by three equal masses
$M_i=\mu_0e^{-1/4}$. These checks do not replace the analytical global
bound. The result rules out ``let a later family fit repair the current
hard radial block''; it does not rule out the Spin(10) action. A
controlled scalar/IR reorganization or a new consistently defined scalar
point must be considered before a physical fit.

\section{A computed linear-\texorpdfstring{$C_{HN}$}{CHN} finite threshold subset}
\subsection{Direct vertex integral and terminal matching}
Use the existing no-half convention
\begin{equation}
 \mathcal L\supset-\tfrac12N^TMN-Y_{\alpha i}(L_\alpha H)N_i
                  +C_{ij}N_iN_j(H^\dagger H)+\mathrm{h.c.},
 \qquad M(\rho)=M-2C\rho,
\end{equation}
where $C=C_{HN}$, $\rho=H^\dagger H$ and $C^T=C$. The Weinberg convention
is $C_{5,\rm tree}^{I}=-YM^{-1}Y^T$. This is the convention used by the
existing RG implementation and its 2024 dimension-five checks
\cite{zhang2024}. Work at leading $q_H/M_i^2=0$ and linear order in $C$.

In a positive Takagi basis, the Yukawa--$C$ bubble contains a
chirality-flipping propagator $M_i/(k^2-M_i^2)$ and a massless scalar
propagator. The no-half $C NN\rho$ vertex supplies a factor two. The
renormalized scalar integral at zero external momenta is
\begin{equation}
 b_i=-\int_0^1\log\frac{tM_i^2}{\mu^2}\,\dd t
       =1+\log\frac{\mu^2}{M_i^2}.
\end{equation}
With the displayed interaction signs, the proper Yukawa correction is
\begin{equation}
 \Delta Y=\frac2{\LOOP}Y\diag(M_i b_i)C.
 \label{eq:CHNY}
\end{equation}
At this order a one-$C$ Higgs tadpole changes its mass but not its
kinetic metric; the corresponding massless Higgs tadpole on a sterile
line is scaleless. This does not discard other canonical type-I wave
functions, which belong to different spurion contributions.

When all sterile fields are removed, the linear-$C$ contribution from
this vertex is
\begin{equation}
 \Delta C_5^{(C)}=-\Delta Y M^{-1}Y^T-YM^{-1}\Delta Y^T.
 \label{eq:CHNC5}
\end{equation}
Differentiating the same fermion determinant at $M(\rho)=M-2C\rho$ gives
the finite Higgs mass contribution
\begin{equation}
 \Delta q_H^{(C)}=\frac4{\LOOP}\sum_i M_i^3\operatorname{Re}C_{ii}
               \left(\log\frac{M_i^2}{\mu^2}-1\right).
 \label{eq:CHNqH}
\end{equation}
Here $q_H$ denotes the coefficient of $+\rho$ in the potential, in the
repository's reference units. Equations~\eqref{eq:CHNY}--\eqref{eq:CHNqH}
are a calculated subset, not the full $C_{HN}$ seesaw matching map.

Their scale derivatives reproduce the independently implemented
$C$-linear beta differences. For example,
\begin{equation}
 \frac{\dd\Delta C_5^{(C)}}{\dd\log\mu}
 =-\frac4{\LOOP}
 \left(YMC M^{-1}Y^T+YM^{-1}CMY^T\right).
 \label{eq:CHNRG}
\end{equation}
The masses in this expression are diagonal positive Takagi masses.
The implementation tests generic complex family covariance, degenerate
$O(3)$ invariance, an independent Feynman-parameter integral and direct
fermion-determinant differentiation. Finite constants are supplied by
the integral, not inferred solely from a beta function.

\subsection{Tree elimination generates a further operator}
Partition the Takagi basis into removed $h$ and retained $r$ blocks.
Although $M_{hr}=0$ at $\rho=0$, generally $C_{hr}\ne0$. Exact algebraic
elimination at nonzero $\rho$ gives
\begin{align}
 Y_r^{\rm eff}(\rho)
 &=Y_r-Y_h(M_h-2C_{hh}\rho)^{-1}(-2C_{hr}\rho)\\
 &=Y_r+D_6\rho+O(\rho^2),\qquad
 D_6=2Y_hM_h^{-1}C_{hr}.
 \label{eq:D6}
\end{align}
In the present sign convention the Lagrangian therefore contains
$-D_6(LH)N_r(H^\dagger H)+\mathrm{h.c.}$.
This is a known class of dimension-six sterile-neutrino operators
\cite{basis}; its coefficient here is derived from the existing action,
not added as a new free matrix. The same expansion gives
\begin{align}
 M_r^{\rm eff}(\rho)&=M_r-2C_{rr}\rho
              -4C_{rh}M_h^{-1}C_{hr}\rho^2+O(\rho^3),\\
 C_5^{\rm eff}(\rho)&=C_5-Y_hM_h^{-1}Y_h^T
              -2Y_hM_h^{-1}C_{hh}M_h^{-1}Y_h^T\rho+O(\rho^2).
\end{align}
Only Eq.~\eqref{eq:D6} is needed for the particular linear-$C$ mixed
threshold consistency check below; the other generated higher operators
have not been declared irrelevant for full matching.

\subsection{A precise mixed-block RG residual}
Let $\mathcal C=C_5-Y_rM_r^{-1}Y_r^T$ denote the active-seesaw composite
below the threshold. In Eq.~\eqref{eq:CHNY} restrict the internal massive
line to the removed block. This determines the hard $\Delta Y_r$ and
the heavy--heavy contribution to $\Delta C_5$. Their contribution to
$\Delta\mathcal C$ is
\begin{equation}
 \Delta\mathcal C_h=\Delta C_{5,hh}
 -\Delta Y_rM_r^{-1}Y_r^T-Y_rM_r^{-1}\Delta Y_r^T.
\end{equation}
Subtracting the existing dimension-five-only beta functions leaves the
identity
\begin{equation}
 \frac{\dd\Delta\mathcal C_h}{\dd\log\mu}
 -\big(\beta_{\mathcal C}^{r}-\beta_{\mathcal C}^{h+r}\big)_{C}
 =\frac2{\LOOP}\left(D_6M_rY_r^T+Y_rM_rD_6^T\right).
 \label{eq:D6mix}
\end{equation}
One can see its origin without fitting coefficients: the omitted terms
in Eq.~\eqref{eq:CHNRG} have an internal retained index and an external
removed index. Replacing their removed propagator by $M_h^{-1}$ gives
the tree vertex in Eq.~\eqref{eq:D6}. A retained $N_r$--Higgs bubble with
one $D_6$ and one $Y_r$ insertion contributes
\begin{equation}
 \beta_{C_5}\big|_{D_6 M_r}=
 \frac2{\LOOP}(D_6M_rY_r^T+Y_rM_rD_6^T)
\end{equation}
in this convention, closing precisely this residual. The finite
retained-line bubble is an EFT graph to subtract, not a missing hard
term to add a second time to $C_5$.

This is not a claim that every dimension-six contribution must always be
kept in a leading dimension-five calculation. If $M_r/M_h$ is treated
as a further small expansion parameter, the displayed effect can be
consistently higher order. However, a finite matching calculation that
keeps the retained masses without this extra hierarchy truncation needs
the generated operator or an explicit bound on what was dropped. No such
bound has yet been established for a physical fitted P54 trajectory.

The actual scalar-tree boundary can have $C\propto M$, but the existing
RG equations do not preserve that alignment for generic $Y_\nu$.
Reading the old two synthetic-family trajectories at their live thresholds
gives $\|D_6\|_F\omega^2$ of
\begin{equation}
 (1.1505\times10^{-4},7.9314\times10^{-5},0),\qquad
 (1.1899\times10^{-4},1.5040\times10^{-4},0).
\end{equation}
The last entry vanishes because no sterile field remains. These are
operator-closure diagnostics, not fitted observables. The old trajectories
were inspected without applying the incomplete new matcher. Full
dimension-six running, nonzero-$q_H$ terms, correlated finite $M_R$,
$\lambda$, wave-function maps and the remaining $C_5$/mixed/box diagrams
are still missing.

\section{Updated acceptance gates and reproducibility}
The common-renormalization verifier passes 46/46 checks; the
$C_{HN}$ closure verifier passes 20/20. The strengthened canonical fit
interface passes 18/18 and the existing upper-vector subset remains
38/38. Passing tests includes correctly detecting a failed physical
stability condition; it is not a claim that all physical gates passed.

The shared fit consumer now additionally requires common parameter/tadpole
consistency, whole-background stability, mass-operator control and
$C_{HN}$ operator closure. Its previous doublet and diagram flags alone
are insufficient. The frozen point is not promoted to a fit. No full
Wilson/box or lower gauge/ghost/Yukawa completion is claimed.

The immediate research choices are:
\begin{enumerate}
\item Keep the action and first test a consistently organized scalar
background with the required Goldstone/IR and momentum/gauge treatment.
The theorem above specifies exactly which previous truncation cannot be
rescued by a family fit; it does not substitute for this further analysis.
\item If a controlled background cannot be obtained at this scalar point,
reselect common scalar parameters subject to stationarity and full radial
and nonradial stability, then recompute the common spectrum/matching.
Changing only the finite counterterm or fitting an isolated mass formula
is not a consistent alternative. No such parameter change is made here.
\item For sequential seesaw, choose and document either the enlarged
operator basis with retained-mass feedback or a justified hierarchy
expansion and remainder bound. Then finish the finite diagrams and
lower matching before the complete light-state and flavor fit.
\end{enumerate}

New artifacts, relative to \texttt{route\_f/}, are
\path{code/verify_p54_common_renormalization.py} and
\path{code/verify_p54_chn_threshold_closure.py}, with matching JSON/Markdown
files under \path{output/}. The first takes \texttt{--cache-dir} pointing
to the old \path{tmp/p54_full_doublet_cw} cache. Only three new mixed
polynomial Hessians were needed, stored in the worktree cache; all later
runs reuse them. Reports record current source hashes, matrices,
individual checks and explicit non-completion flags.

\begin{thebibliography}{9}
\bibitem{vev} V. D\=ud\.enas and M. L\"oschner,
Vacuum expectation value renormalization in the Standard Model and beyond,
\href{https://arxiv.org/abs/2010.15076}{arXiv:2010.15076}.
\bibitem{matching} J. Braathen, M. D. Goodsell and P. Slavich,
Matching renormalisable couplings: simple schemes and a plot,
especially Secs.~2.3 and 4,
\href{https://arxiv.org/html/1810.09388v2}{arXiv:1810.09388v2}.
\bibitem{zhang2024} D. Zhang,
Threshold Effects on the Massless Neutrino in the Canonical Seesaw Mechanism,
\href{https://arxiv.org/html/2405.18017v2}{arXiv:2405.18017v2}.
\bibitem{basis} H.-L. Li, Z. Ren, M.-L. Xiao, J.-H. Yu and Y.-H. Zheng,
Operator Bases in Effective Field Theories with Sterile Neutrinos:
$d\leq9$,
\href{https://arxiv.org/abs/2105.09329}{arXiv:2105.09329}.
\end{thebibliography}
\end{document}
